\documentclass[11pt]{report}
\usepackage[a4paper,margin=2.5cm]{geometry}
\usepackage[T1]{fontenc}
\usepackage{lmodern}
\usepackage{amsmath,amssymb,amsthm,mathtools}
\usepackage{booktabs,array,enumitem}
\usepackage{xcolor}
\usepackage{tikz}
\usetikzlibrary{arrows.meta,positioning,calc,fit,matrix,decorations.pathreplacing}
\usepackage[colorlinks=true,linkcolor=blue!50!black,citecolor=blue!50!black,urlcolor=blue!50!black]{hyperref}
\hypersetup{
  pdftitle={The Spacetime--Electromagnetic Framework, Part II. Field Dynamics, Clifford Transport and a One-Family Pati--Salam Construction},
  pdfauthor={Driss Bencheikh},
  pdfsubject={Clifford transport and a one-family Pati--Salam construction},
  pdfcreator={pdfLaTeX}}
\usepackage{cleveref}
\emergencystretch=1.5em
\setlength{\parskip}{0pt}

% ---------------------------------------------------------------- theorem environments (as in Part I)
\theoremstyle{plain}
\newtheorem{theorem}{Theorem}[chapter]
\newcommand{\sharedthm}[2]{\newtheorem{#1}[theorem]{#2}\crefname{#1}{#2}{#2s}\AddToHook{env/#1/begin}{\crefalias{theorem}{#1}}}
\sharedthm{proposition}{Proposition}
\sharedthm{lemma}{Lemma}
\sharedthm{corollary}{Corollary}
\theoremstyle{definition}
\sharedthm{definition}{Definition}
\sharedthm{hypothesis}{Hypothesis}
\sharedthm{example}{Example}
\sharedthm{remark}{Remark}
\crefname{theorem}{Theorem}{Theorems}
\crefname{chapter}{Chapter}{Chapters}
\crefname{section}{Section}{Sections}
\crefname{appendix}{Appendix}{Appendices}
\crefname{table}{Table}{Tables}
\crefname{figure}{Figure}{Figures}
\crefname{equation}{}{}
\creflabelformat{equation}{(#2#1#3)}
\numberwithin{equation}{chapter}

% ---------------------------------------------------------------- notation (as in Part I)
\newcommand{\R}{\mathbb{R}}
\newcommand{\C}{\mathbb{C}}
\newcommand{\Cl}{\mathrm{Cl}}
\newcommand{\so}[1]{\mathfrak{so}(#1)}
\newcommand{\su}[1]{\mathfrak{su}(#1)}
\newcommand{\uu}[1]{\mathfrak{u}(#1)}
\newcommand{\spp}[1]{\mathfrak{sp}(#1)}
\newcommand{\Tr}{\operatorname{Tr}}
\newcommand{\id}{\mathbf 1}
\renewcommand{\Re}{\operatorname{Re}}
\renewcommand{\Im}{\operatorname{Im}}
\newcommand{\lG}{\mathfrak l_G}
\newcommand{\lint}{\mathfrak l_{\rm int}}
\newcommand{\bB}{\mathbf B}
\newcommand{\bE}{\mathbf E}
\newcommand{\bj}{\mathbf j}
\newcommand{\bk}{\mathbf k}
\newcommand{\bp}{\mathbf p}
\newcommand{\bC}{\mathbf C}
\newcommand{\ba}{\mathbf a}
\newcommand{\Btext}{\mathbf B_{\rm t}}
\newcommand{\half}{\tfrac12}
\newcommand{\curl}{\nabla\times}
\newcommand{\dive}{\nabla\cdot}
\newcommand{\clock}{\hat n}
\newcommand{\gn}{\mathfrak g_n}
\newcommand{\kn}{\mathfrak k_n}
\newcommand{\cB}{\mathcal B}
\newcommand{\cE}{\mathcal E}
\newcommand{\cN}{\mathcal N}

% ---------------------------------------------------------------- Part II conventions
% \PI{Thm~4.12}: a statement of Part I, cited by its printed number.
% \checks{ids}: closes a statement or its proof with the scripted checks that test it (Appendix A), as \checks does in Part I.
% \chk{name}: a script or an identifier that the text itself discusses.
\newcommand{\PI}[1]{Part~I, #1}
\newcommand{\checks}[1]{\par\smallskip{\small\textit{Checks:} #1 (\cref{app:verification}).}}
\newcommand{\chk}[1]{\texttt{#1}}
% \codearchive: where the verification code is published; at a release, replace by its DOI or URL.
\newcommand{\codearchive}{the archive that accompanies this manuscript}
